\section*{Introduction}

The characterization of human population structure---how genetic variation is partitioned among geographic and ethnic groups---is foundational to human genetics. Population structure informs the design of genome-wide association studies \citep{price2006, marchini2010}, guides forensic identification \citep{phillips2015}, shapes our understanding of human evolutionary history \citep{cavalli1994}, and constrains the interpretation of clinical genetic variants across ancestry groups \citep{martin2019}. Over the past two decades, increasingly comprehensive catalogs of human genetic variation have refined our understanding of population differentiation, from early microsatellite surveys \citep{rosenberg2002} to dense SNP arrays \citep{li2008} and whole-genome sequencing panels \citep{1000genomes2015, byrskaBishop2022}.

Yet this body of work has been built almost entirely on a single class of genetic variation: single nucleotide variants. SNVs dominate population genetic analyses for practical reasons---they are abundant, well-genotyped, and amenable to standard analytical tools---but they represent only one facet of the mutational spectrum. Human genomes also harbor millions of short insertions and deletions (INDELs, typically 1--50~bp) and tens of thousands of structural variants (SVs, typically $\geq$50~bp), including deletions, insertions, duplications, and inversions \citep{sudmant2015, collins2020, byrskaBishop2022}. These larger variant classes differ fundamentally from SNVs in their mutational mechanisms, functional impact, and selective regimes: SVs have higher per-event mutation rates, disrupt more sequence per event, and are subject to stronger purifying selection \citep{abel2020, audano2019}. Whether these differences produce concordant or discordant population-genetic signals remains largely unexplored.

The question is not merely academic. If different variant types encode the same population structure, then SNV-centric analyses faithfully represent the genome-wide picture, and the field's methodological choices are validated. If, however, variant types diverge in their population-genetic signals, then relying exclusively on SNVs risks missing variant-type-specific evolutionary information---and tools designed for one variant class may not generalize to others. This distinction has practical implications for ancestry inference, selection scans, and clinical interpretation of non-SNV variants.

Theoretical considerations suggest that concordance and discordance should coexist. To the extent that population structure reflects shared demographic history---migration, drift, and admixture---all variant types should recover the same topology, because these forces act genome-wide regardless of the mutational mechanism that generated a variant \citep{wright1949}. However, differential selection across variant classes could produce systematic quantitative discordance: if SVs are subject to stronger purifying selection than SNVs, their allele frequencies should be shifted toward rare variants, reducing their contribution to inter-population differentiation and attenuating $F_\mathrm{ST}$ \citep{charlesworth1998}. The expected signature is concordant topology but discordant magnitude---the same populations should cluster together, but the distances between them should differ depending on variant type.

Individual variant classes have been studied in population-genetic contexts---\citet{sudmant2015} characterized SV-based population structure and $F_\mathrm{ST}$ patterns in the Phase 3 1000 Genomes panel, while \citet{collins2020} expanded the SV catalog---but no study has systematically quantified the concordance of population structure, differentiation magnitude, and allele frequency dynamics \emph{across} SNVs, INDELs, and SVs within a unified analytical framework. In particular, the contribution of variant count versus biology to observed cross-type discordance, the relationship between variant size and differentiation, and the question of whether insertions and deletions behave symmetrically have not been addressed. A recent comparison of SV, INDEL, and SNP differentiation in locally adapted Atlantic salmon populations \citep{lecomte2024} found that all three variant types contributed to population differentiation but did not examine size-dependent effects or selective asymmetries in humans. In humans, the 2022 high-coverage release of the 1000 Genomes Project \citep{byrskaBishop2022}---with uniformly genotyped SNVs, INDELs, and SVs across 3,202 individuals from 26 globally distributed populations---provides an unprecedented opportunity to conduct this comparison.

Here we exploit this resource to ask: do different variant types tell the same population-genetic story, and if not, where and why do they diverge? We quantify concordance using Procrustes analysis of PCA embeddings, $F_\mathrm{ST}$ correlation across 325 population pairs, and site frequency spectrum comparisons. We develop a convergence analysis framework that disentangles statistical power effects from genuine biological discordance by systematically varying the number of variants used for PCA. We further test whether differentiation scales with variant size by computing $F_\mathrm{ST}$ across size-stratified bins spanning two orders of magnitude, and we decompose SVs into insertions and deletions to test for asymmetric selective constraint.

Our analyses reveal three principal findings. First, while all variant types recover the canonical continental population structure, SVs encode a quantitatively attenuated differentiation signal that plateaus at a Procrustes concordance of $\sim$0.60 with SNVs regardless of the number of markers---establishing the discordance as irreducible biology, not a statistical artifact. Second, population differentiation decreases monotonically with variant size along a continuous gradient from 1~bp INDELs to $>$50~bp SVs, consistent with stronger purifying selection on larger variants. Third, insertions show 25\% lower $F_\mathrm{ST}$ than deletions of comparable size, suggesting asymmetric selective constraint that challenges the common assumption of symmetric insertion-deletion dynamics. Together, these findings establish a size-dependent selection gradient as a previously unrecognized axis of human population genomic variation.
